\begin{abstract}
La fabricación aditiva (AM) representa una tecnología disruptiva para la industria satelital de telecomunicaciones, permitiendo la producción de componentes ligeros, complejos y funcionales integrados que reducen significativamente el peso y los costos de lanzamiento. Este artículo revisa el estado del arte en AM aplicada a estructuras de plataformas, antenas RF y subsistemas, destacando reducciones de masa de hasta 30-40% y consolidación de piezas que optimizan ensamblaje y mantenimiento. Se presenta una metodología de rediseño DfAM para plataformas minisatélite con estructuras lattice, junto con análisis de desempeño RF en bandas Ku/Ka. Los resultados demuestran cumplimiento de requisitos dinámicos y térmicos con menor lead-time y costos. Se discuten desafíos de calificación espacial y se proponen direcciones futuras para manufactura in-situ. Los hallazgos apoyan la adopción de AM para constelaciones de satélites de bajo costo y alto rendimiento en telecomunicaciones.
\end{abstract}